\chapter{نخستین برنامه}
\label{ch:first}

\begin{goals}
\begin{itemize}
  \item برنامه چیست و رایانه چگونه آن را اجرا می‌کند
  \item آشنایی با ویرایشگر برخط سلام و اجرای نخستین برنامه
  \item چاپ کردن متن و عدد روی صفحه با \fcode{سرچاپ} و \fcode{چاپ}
  \item نوشتن توضیح در میان برنامه
  \item خواندن پیام خطا و رفع ساده‌ترین اشتباه‌ها
\end{itemize}
\end{goals}

\section{برنامه چیست؟}

فرض کنید می‌خواهید به دوستی که هرگز چای دم نکرده است، یاد بدهید چای دم
کند. نمی‌توانید بگویید «یک چای خوب درست کن»؛ باید قدم‌به‌قدم بگویید: کتری
را پر از آب کن، روی اجاق بگذار، صبر کن تا بجوشد، چای خشک را در قوری بریز
و همین‌طور تا آخر. رایانه هم درست چنین دوستی است: بسیار سریع، بسیار دقیق
و بسیار کم‌هوش. هر کاری را که به او بگویید انجام می‌دهد، اما فقط به شرطی که
دقیقاً بگویید چه کند.

\textbf{برنامه} همین فهرست دستورهای قدم‌به‌قدم است، و \textbf{برنامه‌نویسی}
هنر نوشتن آن. زبانی که این دستورها را با آن می‌نویسیم
\textbf{زبان برنامه‌نویسی} نام دارد. زبان برنامه‌نویسی، برخلاف زبان انسان‌ها،
واژه‌های اندکی دارد و قاعده‌هایش بسیار سخت‌گیرانه است؛ اما در عوض هیچ
ابهامی در آن نیست. زبان این کتاب، \textbf{سلام} است: زبانی که دستورهایش به
فارسی نوشته می‌شود.

\begin{note}
وقتی برنامه‌ای می‌نویسید، متنی که نوشته‌اید را \textbf{کد} یا
\textbf{کد منبع} می‌نامند. رایانه این متن را می‌خواند، دستورهایش را یکی‌یکی
انجام می‌دهد و نتیجه را نشان می‌دهد. به این کار \textbf{اجرا کردن} برنامه
می‌گویند.
\end{note}

\section{ویرایشگر برخط سلام}

برای نوشتن و اجرا کردن برنامه‌های این کتاب به هیچ نصب و تنظیمی نیاز
ندارید. مرورگر خود را باز کنید و به این نشانی بروید:
\begin{center}
\url{https://ediotr.salamlang.ir}
\end{center}
صفحه‌ای که می‌بینید «زمین بازی سلام» نام دارد و از سه بخش اصلی ساخته شده
است:
\begin{enumerate}
  \item \textbf{نوار بالا:} در این نوار می‌توانید زبان را انتخاب کنید. گزینه‌ی
        \textbf{فارسی} را برگزینید تا هم ظاهر صفحه فارسی شود و هم ویرایشگر
        آماده‌ی دریافت کد فارسی. در همین نوار دو حالت \textbf{برنامه} و
        \textbf{چیدمان} وجود دارد؛ ما در سراسر این کتاب در حالت
        \textbf{برنامه} کار می‌کنیم. دکمه‌ی \textbf{اجرا} هم همین‌جاست.
  \item \textbf{ویرایشگر:} پنجره‌ای که کد را در آن می‌نویسید. شماره‌ی خط‌ها
        در کنار آن دیده می‌شود و واژه‌های مهم زبان رنگی می‌شوند.
  \item \textbf{خروجی:} پنجره‌ای که نتیجه‌ی اجرای برنامه در آن نمایش داده
        می‌شود. اگر برنامه اشتباهی داشته باشد، پیام خطا هم همین‌جا ظاهر
        می‌شود.
\end{enumerate}

\begin{tip}
در نوار بالای ویرایشگر فهرستی از برنامه‌های نمونه هم هست. بد نیست همین حالا
یکی از آن‌ها را باز کنید و دکمه‌ی اجرا را بزنید تا ببینید همه چیز درست کار
می‌کند. نگران نباشید اگر هنوز چیزی از کد سر در نمی‌آورید؛ تا پایان این
کتاب همه‌ی آن‌ها را خواهید فهمید.
\end{tip}

\section{سلام، دنیا!}

سنتی قدیمی میان برنامه‌نویسان هست: نخستین برنامه در هر زبان تازه، برنامه‌ای
است که جمله‌ی «سلام، دنیا!» را روی صفحه می‌نویسد. ما هم همین سنت را نگه
می‌داریم. متن ویرایشگر را پاک کنید و این سه خط را در آن بنویسید:

\program{سلام، دنیا!}
\begin{salamfa}
روال ریشه:
    سرچاپ "سلام، دنیا!"
پایان
\end{salamfa}

حالا دکمه‌ی \textbf{اجرا} را بزنید. در پنجره‌ی خروجی این را می‌بینید:

\begin{outputfa}
سلام، دنیا!
\end{outputfa}

تبریک! شما همین حالا نخستین برنامه‌ی خود را نوشتید و اجرا کردید. حالا
ببینیم این سه خط چه می‌گویند.

\subsection{روال ریشه: نقطه‌ی شروع}

خط نخست، \fcode{روال ریشه:}، به سلام می‌گوید: «اجرای برنامه از این‌جا شروع
می‌شود.» هر برنامه‌ی سلام باید دقیقاً یک \fcode{روال ریشه} داشته باشد.
واژه‌ی \fcode{پایان} در خط آخر پایان آن را نشان می‌دهد. فعلاً این دو خط را
مانند جلد یک دفتر ببینید: هر چه می‌نویسیم، میان این دو خط می‌نویسیم. در فصل
\ref{ch:functions} خواهیم دید که واژه‌ی «روال» چیزهای بیشتری در چنته دارد.

به دونقطه‌ی پایان خط نخست دقت کنید. در سلام هر جا که یک بلوک (یک دسته
دستور که با هم هستند) آغاز می‌شود، خط با \code{:} تمام می‌شود و بلوک با
\fcode{پایان} بسته می‌شود. این الگو را بارها و بارها خواهید دید.

\subsection{تورفتگی: نظم دیداری برنامه}

دقت کنید که خط میانی کمی جلوتر از دو خط دیگر نوشته شده است. به این
جلوآمدگی \textbf{تورفتگی} می‌گویند و نشان می‌دهد این خط \emph{درون} روال ریشه است. رسم بر این است که هر چه درون یک بلوک است را با چهار فاصله تو
ببریم. سلام برای اجرای برنامه به این فاصله‌ها نیازی ندارد و اگر آن‌ها را
ننویسید هم برنامه اجرا می‌شود؛ اما برنامه‌ای که تورفتگی درستی ندارد، مثل
کتابی است که پاراگراف‌بندی ندارد: می‌شود خواند، اما سخت.

\subsection{چاپ: نوشتن روی صفحه}

خط میانی همان جایی است که کار انجام می‌شود. دستور \kwd{سرچاپ} هر چه را که
جلویش بنویسید روی پنجره‌ی خروجی می‌نویسد و سپس به خط بعد می‌رود. متنی که
می‌خواهیم چاپ شود را میان دو علامت نقل قول \code{"} می‌گذاریم. به این متن
داخل نقل قول \textbf{رشته} می‌گویند؛ در فصل \ref{ch:strings} بسیار با آن
سروکار خواهیم داشت. خود علامت‌های نقل قول چاپ نمی‌شوند؛ آن‌ها فقط به سلام
می‌گویند متن از کجا شروع می‌شود و کجا پایان می‌یابد.

\begin{tip}
علامت استاندارد نقل قول \code{"} است که در صفحه‌کلید معمولاً با نگه‌داشتن
شیفت و زدن کلید \code{'} نوشته می‌شود. اگر تایپ آن در متن فارسی دست‌وپاگیر
بود، می‌توانید به‌جای آن از گیومه‌ی فارسی \code{«...»} هم استفاده کنید:
\fcode{سرچاپ «سلام، دنیا!»} درست همان کاری را می‌کند که
\fcode{سرچاپ "سلام، دنیا!"} می‌کند. با این‌حال \code{"} رایج‌تر و استاندارد
است و در این کتاب همه‌جا همان را به کار می‌بریم. البته می‌توانید \emph{درون}
رشته هر علامتی، از جمله گیومه‌ی فارسی، بنویسید: \fcode{سرچاپ "او گفت: «سلام»"}
کاملاً درست است.
\end{tip}

\section{چند دستور پشت سر هم}

یک برنامه به ندرت فقط یک دستور دارد. دستورها را زیر هم می‌نویسیم و سلام
آن‌ها را به همان ترتیب، از بالا به پایین، اجرا می‌کند:

\program{چند خط چاپ}
\begin{salamfa}
روال ریشه:
    سرچاپ "به کتاب آموزش سلام خوش آمدید."
    سرچاپ "این دومین خط است."
    سرچاپ "و این سومین خط."
پایان
\end{salamfa}

\begin{outputfa}
به کتاب آموزش سلام خوش آمدید.
این دومین خط است.
و این سومین خط.
\end{outputfa}

هر \fcode{سرچاپ} یک خط تازه در خروجی می‌سازد. ترتیب مهم است: اگر جای دو دستور
را عوض کنید، جای دو خط در خروجی هم عوض می‌شود. همین حالا امتحان کنید.

\section{چاپ عددها و محاسبه}

\fcode{سرچاپ} فقط برای متن نیست؛ عدد را هم چاپ می‌کند. عددها را بدون نقل قول
می‌نویسیم:

\begin{salamfa}
روال ریشه:
    سرچاپ ۱۴۰۴
    سرچاپ ۲ + ۳
    سرچاپ ۱۰ * ۴
پایان
\end{salamfa}

\begin{outputfa}
1404
5
40
\end{outputfa}

به خط دوم دقت کنید: نوشتیم \code{۲ + ۳} اما سلام \code{۵} را چاپ کرد. یعنی
سلام پیش از چاپ، حساب را انجام داد. علامت \code{*} در برنامه‌نویسی همان ضرب
است (چون علامت × روی صفحه‌کلید نیست). در فصل \ref{ch:operators} به طور کامل
درباره‌ی محاسبه صحبت می‌کنیم.

حالا ببینید اگر همان \code{۲ + ۳} را داخل نقل قول بگذاریم چه می‌شود:

\begin{salamfa}
روال ریشه:
    سرچاپ "2 + 3"
    سرچاپ ۲ + ۳
پایان
\end{salamfa}

\begin{outputfa}
2 + 3
5
\end{outputfa}

هر چه داخل نقل قول باشد \emph{متن} است و سلام به محتوایش کاری ندارد؛ آن را
همان‌طور که هست چاپ می‌کند. اما بیرون نقل قول، \code{۲ + ۳} یک محاسبه است.
این تفاوت میان «متن» و «عدد» یکی از مهم‌ترین چیزهایی است که در این کتاب یاد
می‌گیرید.

\section{چاپ چند چیز در یک خط}

اگر بخواهیم در یک خط هم متن و هم عدد چاپ کنیم، آن‌ها را با ویرگول
انگلیسی \code{,} از هم جدا می‌کنیم. سلام آن‌ها را پشت سر هم، با یک فاصله
میانشان، چاپ می‌کند:

\program{متن و عدد در یک خط}
\begin{salamfa}
روال ریشه:
    سرچاپ "حاصل جمع ۲ و ۳ برابر است با", ۲ + ۳
    سرچاپ "سن سارا:", ۱۷, "سال"
پایان
\end{salamfa}

\begin{outputfa}
حاصل جمع ۲ و ۳ برابر است با 5
سن سارا: 17 سال
\end{outputfa}

\begin{note}
درون رشته‌ها می‌توانید از رقم‌های فارسی استفاده کنید و همان‌طور که هست
چاپ می‌شوند. اما خود سلام نتیجه‌ی محاسبه‌ها را همیشه با رقم‌های انگلیسی
چاپ می‌کند، حتی اگر عددهای کد را با رقم فارسی نوشته باشید؛ در خروجی
بالا هم به همین دلیل «۵» به شکل انگلیسی \code{5} چاپ شده است. سلام برای
نوشتن یک عدد در کد هر دو رقم را می‌پذیرد: \fcode{چاپ \lr{\salamcodefont ۲ + ۳}} و
\fcode{چاپ \lr{\salamcodefont 2 + 3}} هر دو همان نتیجه را می‌دهند. در این کتاب عددهای کد را
با رقم فارسی می‌نویسیم، چون کتاب فارسی است؛ فقط یادتان باشد که خروجی
خودِ سلام همیشه انگلیسی می‌ماند.
\end{note}

\section{بنویس: چاپ بدون رفتن به خط بعد}

دستور \kwd{چاپ} درست مثل \fcode{سرچاپ} است، با یک تفاوت: بعد از نوشتن، به
خط بعد نمی‌رود. یعنی دستور بعدی از همان‌جا ادامه می‌دهد:

\program{تفاوت بنویس و چاپ}
\begin{salamfa}
روال ریشه:
    چاپ "الف"
    چاپ "ب"
    چاپ "پ"
    سرچاپ ""
    سرچاپ "ت"
    سرچاپ "ث"
پایان
\end{salamfa}

\begin{outputfa}
الفبپ
ت
ث
\end{outputfa}

سه \fcode{چاپ} نخست همه در یک خط نوشتند. سپس \fcode{سرچاپ ""} یک رشته‌ی
خالی چاپ کرد؛ یعنی چیزی ننوشت اما به خط بعد رفت. \fcode{چاپ} در فصل‌های
بعد، وقتی بخواهیم چیزی را قطعه‌قطعه در یک خط بسازیم، به کارمان می‌آید. تا
آن زمان بیشتر از \fcode{سرچاپ} استفاده می‌کنیم.

\begin{tip}
\fcode{سرچاپ} را می‌توانید بدون هیچ چیز هم بنویسید. \fcode{سرچاپ} به تنهایی
یک خط خالی در خروجی می‌سازد و برای فاصله انداختن میان بخش‌های خروجی مفید
است.
\end{tip}

\section{توضیح نوشتن در برنامه}

گاهی می‌خواهیم در میان کد، برای خودمان یا دیگران یادداشتی بگذاریم:
این خط چه کار می‌کند، چرا این‌طور نوشته شده، یا چه چیزی هنوز کامل نیست.
به این یادداشت‌ها \textbf{توضیح} می‌گویند. سلام هر چه را بعد از دو خط مورب
\code{//} بیاید تا پایان همان خط نادیده می‌گیرد:

\program{توضیح در برنامه}
\begin{salamfa}
// این برنامه یک پیام خوش‌آمد چاپ می‌کند
روال ریشه:
    سرچاپ "خوش آمدید"    // چاپ پیام اصلی
    // چاپ "این خط اجرا نمی‌شود"
    سرچاپ "روز خوبی داشته باشید"
پایان
\end{salamfa}

\begin{outputfa}
خوش آمدید
روز خوبی داشته باشید
\end{outputfa}

خط چهارم را ببینید: با اینکه یک دستور \fcode{سرچاپ} است، چون جلویش \code{//}
گذاشته‌ایم اجرا نشده است. برنامه‌نویسان از این ترفند برای «خاموش کردن»
موقت یک خط استفاده می‌کنند، بی‌آنکه آن را پاک کنند.

اگر توضیح چند خط باشد، می‌توانید آن را میان \code{/*} و \code{*/} بنویسید:

\begin{salamfa}
/*
   این برنامه را سارا نوشته است.
   کار آن: چاپ یک شعر کوتاه.
*/
روال ریشه:
    سرچاپ "بنی آدم اعضای یکدیگرند"
    سرچاپ "که در آفرینش ز یک گوهرند"
پایان
\end{salamfa}

\begin{outputfa}
بنی آدم اعضای یکدیگرند
که در آفرینش ز یک گوهرند
\end{outputfa}

\begin{tip}
توضیح خوب چیزی را می‌گوید که از خود کد پیدا نیست: «چرا» را می‌گوید، نه
«چه» را. نوشتن \fcode{// چاپ پیام} بالای \fcode{سرچاپ "پیام"} کمکی به کسی
نمی‌کند؛ اما نوشتن \fcode{// این عدد را از جدول قیمت‌های ۱۴۰۴ برداشته‌ایم}
ارزشمند است. در این کتاب، برای آموزش، بیشتر از حد معمول توضیح می‌نویسیم.
\end{tip}

\section{وقتی اشتباه می‌کنیم}

همه اشتباه می‌کنند؛ برنامه‌نویسان باتجربه فقط سریع‌تر اشتباهشان را پیدا
می‌کنند. بیایید عمداً اشتباه کنیم تا ببینیم سلام چه می‌گوید. علامت نقل قول
پایانی را جا می‌اندازیم:

\begin{salamfa}
روال ریشه:
    سرچاپ "سلام، دنیا!
پایان
\end{salamfa}

با زدن دکمه‌ی اجرا، به جای خروجی، پیامی مانند این می‌بینید:

\begin{outputfa}
خطا: رشته‌ی پایان‌نیافته
 --> 2:10
\end{outputfa}

پیام خطا دو چیز مهم به ما می‌گوید: \textbf{چه} اشتباهی رخ داده («رشته‌ی
پایان‌نیافته»، یعنی رشته‌ای که نقل قول پایانی ندارد) و \textbf{کجا}
(خط ۲). عادت کنید پیام‌های خطا را با دقت بخوانید؛ آن‌ها دشمن شما نیستند،
راهنمای شما هستند. در پیوست \ref{app:errors} فهرستی از خطاهای رایج و معنی
آن‌ها را آورده‌ایم.

چند اشتباه بسیار رایج در همین نخستین قدم‌ها:
\begin{itemize}
  \item فراموش کردن \fcode{پایان} در پایان روال ریشه.
  \item فراموش کردن دونقطه‌ی بعد از \fcode{روال ریشه}.
  \item فراموش کردن ویرگول میان چیزهایی که با یک \fcode{سرچاپ} چاپ می‌شوند.
  \item نوشتن نام دستورها با املای اشتباه، مثلاً \fcode{چاپچ} به جای
        \fcode{سرچاپ}.
\end{itemize}

\begin{deep}
چگونه رایانه‌ای که فقط صفر و یک می‌فهمد، متن فارسی ما را اجرا می‌کند؟
پاسخ برنامه‌ای است به نام \textbf{مترجم} (یا کامپایلر). مترجم سلام کد شما
را می‌خواند، اول بررسی می‌کند که قاعده‌های زبان را رعایت کرده باشد (همان‌جا
که خطاها را می‌بینید) و سپس آن را به دستورهایی تبدیل می‌کند که پردازنده‌ی
رایانه مستقیم می‌فهمد. در ویرایشگر برخط، همین مترجم درون مرورگر شما اجرا
می‌شود. جالب است بدانید که مترجم سلام، خودش با زبان سلام نوشته شده است.
\end{deep}

\section{یک برنامه‌ی کامل‌تر}

بیایید هر چه در این فصل آموختیم را در یک برنامه کنار هم بگذاریم: کارت
شناسایی کوچکی که خودتان را معرفی می‌کند.

\program{کارت معرفی}
\begin{salamfa}
// کارت معرفی
روال ریشه:
    سرچاپ "================="
    سرچاپ "نام: سارا احمدی"
    سرچاپ "شهر: کاشان"
    سرچاپ "سن:", ۲۰۲۶ - ۲۰۰۸, "سال"
    سرچاپ "علاقه‌مندی: برنامه‌نویسی"
    سرچاپ "================="
پایان
\end{salamfa}

\begin{outputfa}
=================
نام: سارا احمدی
شهر: کاشان
سن: 18 سال
علاقه‌مندی: برنامه‌نویسی
=================
\end{outputfa}

خط «سن» را ببینید: به جای نوشتن عدد ۱۸، تفریق دو سال را نوشتیم و سلام خودش
حساب کرد. این یک عادت خوب است: هر جا عددی از محاسبه به دست می‌آید، بگذارید
رایانه حساب کند، نه شما.

\begin{summary}
\begin{itemize}
  \item برنامه فهرستی از دستورهاست که رایانه به ترتیب، از بالا به پایین،
        اجرا می‌کند.
  \item هر برنامه‌ی سلام با \fcode{روال ریشه:} شروع و با \fcode{پایان}
        پایان می‌یابد. دستورها میان این دو خط و با تورفتگی نوشته می‌شوند.
  \item \fcode{سرچاپ} چیزی را می‌نویسد و به خط بعد می‌رود؛ \fcode{چاپ}
        می‌نویسد اما در همان خط می‌ماند.
  \item متن را میان \code{"} می‌نویسیم و عدد را بدون آن. چند چیز را با
        \code{,} از هم جدا می‌کنیم.
  \item هر چه بعد از \code{//} بیاید توضیح است و اجرا نمی‌شود.
  \item پیام خطا می‌گوید چه اشتباهی در کدام خط رخ داده است. آن را با دقت
        بخوانید.
\end{itemize}
\end{summary}

\section*{تمرین‌ها}

\exercise برنامه‌ای بنویسید که نام، نام خانوادگی و شهر شما را هر یک در یک
خط جداگانه چاپ کند.

\exercise برنامه‌ای بنویسید که با سه دستور \fcode{چاپ} و یک دستور
\fcode{سرچاپ}، جمله‌ی «من برنامه‌نویس هستم» را در یک خط چاپ کند.

\exercise برنامه‌ای بنویسید که این محاسبه‌ها را انجام دهد و نتیجه‌ی هر یک
را در یک خط، همراه با توضیحی کوتاه، چاپ کند: حاصل جمع ۱۲۵ و ۳۷۵؛
حاصل ضرب ۲۴ در ۷؛ تعداد ساعت‌های یک هفته.

\exercise برنامه‌ی زیر سه اشتباه دارد. بدون اجرا کردن، آن‌ها را پیدا کنید و
سپس با اجرا کردن پاسخ خود را بررسی کنید:
\begin{salamfa}
روال ریشه
    print "سلام"
    چاپ "خداحافظ
پایان
\end{salamfa}

\exercise برنامه‌ای بنویسید که یک شعر یا ضرب‌المثل چهارخطی دلخواه را چاپ
کند و بالای برنامه، در یک توضیح چندخطی، نام شاعر یا منبع آن را بنویسید.
